% Standalone problem document, generated from the adjacent README.md.
% Compile with XeLaTeX. Mathematical content is maintained in Markdown.
\documentclass[11pt,a4paper]{article}
\usepackage[fontsize=10.5pt]{scrextend}
\usepackage[margin=22mm,headheight=15pt,headsep=9mm,footskip=11mm]{geometry}
\usepackage{fontspec}
\setmainfont{texgyrepagella}[Extension=.otf,UprightFont=*-regular,BoldFont=*-bold,ItalicFont=*-italic,BoldItalicFont=*-bolditalic]
\setsansfont{texgyreheros}[Extension=.otf,UprightFont=*-regular,BoldFont=*-bold,ItalicFont=*-italic,BoldItalicFont=*-bolditalic]
\setmonofont{texgyrecursor}[Extension=.otf,UprightFont=*-regular,BoldFont=*-bold,ItalicFont=*-italic,BoldItalicFont=*-bolditalic,Scale=MatchLowercase]
\usepackage{amsmath,amssymb}
\usepackage{unicode-math}
\setmathfont{texgyrepagella-math.otf}
\usepackage{xcolor,microtype,enumitem,needspace,fancyhdr}
\usepackage{bookmark}
\definecolor{ink}{HTML}{17344B}
\definecolor{accent}{HTML}{197D80}
\definecolor{muted}{HTML}{596773}
\hypersetup{colorlinks=true,linkcolor=accent,urlcolor=accent,pdftitle={IE-16 - A
sharp subset bound for worst-case normal
GMRES},pdfauthor={Open Problems in Numerical Linear Algebra}}
\urlstyle{same}
\setlength{\parindent}{0pt}
\setlength{\parskip}{5pt plus 1pt minus 1pt}
\linespread{1.04}
\setlength{\emergencystretch}{3em}
\widowpenalty=10000
\clubpenalty=10000
\displaywidowpenalty=10000
\allowdisplaybreaks[1]
\setlist{nosep,leftmargin=1.5em}
\providecommand{\tightlist}{\setlength{\itemsep}{0pt}\setlength{\parskip}{0pt}}
\providecommand{\pandocbounded}[1]{#1}
\pagestyle{fancy}
\fancyhf{}
\fancyhead[L]{\sffamily\footnotesize\color{muted}OPEN PROBLEMS IN NUMERICAL LINEAR ALGEBRA}
\fancyhead[R]{\sffamily\bfseries\footnotesize\color{accent}IE-16}
\fancyfoot[L]{\parbox[b]{0.9\textwidth}{\sffamily\fontsize{8}{10}\selectfont\color{muted}Editorial ratings; a bounded search cannot certify that a problem remains unsolved.\\Literature check: 2026-09-12}}
\fancyfoot[R]{\sffamily\footnotesize\color{muted}\thepage}
\renewcommand{\headrulewidth}{0.3pt}
\renewcommand{\headrule}{\hbox to\headwidth{\color{accent}\leaders\hrule height \headrulewidth\hfill}}
\makeatletter
\renewcommand{\section}{\Needspace{5\baselineskip}\@startsection{section}{1}{0pt}{12pt plus 2pt minus 2pt}{4pt}{\sffamily\bfseries\large\color{ink}}}
\renewcommand{\subsection}{\Needspace{4\baselineskip}\@startsection{subsection}{2}{0pt}{9pt plus 1pt minus 1pt}{3pt}{\sffamily\bfseries\normalsize\color{ink}}}
\makeatother
\setcounter{secnumdepth}{0}
\begin{document}
{\sffamily\small\color{muted}Linear systems and elimination\par}
\vspace{3pt}
{\raggedright\hyphenpenalty=10000\exhyphenpenalty=10000\sffamily\bfseries\fontsize{21}{24}\selectfont\color{ink}A
sharp subset bound for worst-case normal GMRES\par}
\vspace{7pt}
{\sffamily\small\textbf{Difficulty:} challenging\quad\textbf{Importance:} interesting
to the community\par}
{\sffamily\small\bfseries\color{accent}Status: Solved\par}
\vspace{4pt}
{\color{accent}\hrule height 0.8pt}
\vspace{6pt}
\textbf{Historical rating rationale:} Challenging reflects extending a
discrete real approximation principle to complex spectra with a sharp
constant; community impact concerns usable worst-case GMRES bounds for
normal matrices.

\subsection{Resolution: negative, 12 September
2026}\label{resolution-negative-12-september-2026}

\textbf{Author:} Sidney Holden, Center for Computational Biology,
Flatiron Institute, Simons Foundation.
\href{https://github.com/ajt60gaibb/OpenProblemsInNLA/blob/main/references/holden-ie16-2026-09-12/README.md}{Verified
affiliation and submission record}.

\href{https://github.com/ajt60gaibb/OpenProblemsInNLA/blob/main/linear-systems-and-elimination/IE-16/solution.pdf}{Theorem
1.1, proved in Sections 2-3 of the manuscript} disproves the displayed
universal inequality. With \(\omega=e^{2\pi i/3}\), take the nine
distinct nonzero points \(L=\{\omega^a+10^{-3}\omega^b:0\le a,b\le2\}\)
and \(k=4\). Then

\[M_4(L)=\frac{3003003000}{1001003001001}>\frac{299}{100000},
\qquad \max_{S\subseteq L,\ |S|=5}M_4(S)<\frac{23}{10000}.\]

The ratio is therefore greater than \(13/10>4/\pi\). Here \(n=9\) and
\(1\le4\le n-2\), so all original assumptions and quantifiers are
addressed by this counterexample. Theorem 1.2 (Sections 4-5)
additionally proves that no finite dimension-independent replacement
constant exists. The finite counterexample alone settles IE-16
negatively; no case of the original universal truth question remains
open.

The complete proof passed a separate
\href{https://github.com/ajt60gaibb/OpenProblemsInNLA/blob/main/references/holden-ie16-2026-09-12/INDEPENDENT-REVIEW.md}{independent
Codex AI-agent informal audit}. The supplied exact-arithmetic verifier
passed all 126 subset checks, and reviewer-written rational checks
passed independently. This is not external human peer review or formal
verification. No Lean verification was performed.
\href{https://github.com/ajt60gaibb/OpenProblemsInNLA/blob/main/linear-systems-and-elimination/IE-16/solution.tex}{Proof
source} ·
\href{https://github.com/ajt60gaibb/OpenProblemsInNLA/blob/main/references/holden-ie16-2026-09-12/submitted/README.md}{Exact
verifier and certificates}.

\subsection{Original problem
(retained)}\label{original-problem-retained}

Let \(L\subset\mathbb C\setminus\{0\}\) consist of \(n\ge3\) distinct
points. For a nonempty \(S\subseteq L\), define

\[M_k(S)=\min_{p\in\mathbb C[z],\ \deg p\le k,\ p(0)=1}\max_{z\in S}|p(z)|.\]

Is it true, simultaneously for every such \(L\) and \(1\le k\le n-2\),
that

\[M_k(L)\le\frac4\pi\max_{S\subseteq L,\ |S|=k+1}M_k(S)?\]

The constant is required to be independent of dimension, degree, and the
locations of the points.

For a normal matrix with spectrum \(L\), \(M_k(L)\) is the largest
relative residual norm attainable at GMRES step \(k\), over unit initial
residuals. The smaller-set quantities have explicit formulas through
Lagrange interpolation. The question therefore asks for a sharp,
dimension-independent certificate of worst-case convergence using small
subsets of the spectrum. It is distinct from equality of ideal and
worst-case GMRES for a nonnormal Jordan block (IE-02).

\subsection{References}\label{references}

J. Liesen and P. Tichý, \emph{The worst-case GMRES for normal matrices},
BIT 44 (2004), 79--98, conjecture (3.16), pp.~91--92, and Appendix
(\href{https://page.math.tu-berlin.de/~liesen/Publicat/LieTic04.pdf}{author
copy}). Their \emph{A min-max problem on roots of unity}, TU Berlin
Preprint 28-2003, conclusion (9.1), proves selected roots-of-unity cases
and sharpness of \(4/\pi\)
(\href{https://doi.org/10.14279/depositonce-14263}{university record}).
Their survey \emph{Convergence analysis of Krylov subspace methods},
GAMM-Mitteilungen 27 (2004), discussion after (11), restates the
conjecture
(\href{https://page.math.tu-berlin.de/~liesen/Publicat/LiTiGAMM.pdf}{author
copy}).

\subsection{Status check --- 2026-09-08}\label{status-check-2026-09-08}

Searched the original title with ``conjecture'', ``4/pi'', ``4/π'',
``proof'', ``counterexample'', and 2025--2026; checked the authors'
current publication lists. The later Jordan-block papers concern a
different conjecture. No resolution or recent explicit reaffirmation of
this particular bound was located; its open-status evidence is
historical and bounded.

\subsection{Audit update --- 2026-09-10}\label{audit-update-2026-09-10}

The
\href{https://page.math.tu-berlin.de/~liesen/Publicat/LieTic04.pdf}{primary
paper}, §3.2.1, proves the stronger subset equality for every real
spectrum, a substantive subclass of the displayed target. Its
complex-spectrum conjecture remains in §3.2.2. Searches for the
Liesen--Tichý conjecture and later sharp normal-GMRES subset bounds
found no general resolution; the current status records the
real-spectrum partial result without asserting exhaustive literature
coverage.
\end{document}
